\subsubsection{ソフトウェア設計の主要論点}

設計では、多くの論点を同時に扱う必要がある。代表的な論点は、性能、セキュリティ、信頼性、使いやすさ、保守容易性などの品質である。また、ソフトウェアをどのような構成要素へ分けるか、それらをどう接続するか、どの単位で責務を持たせるかも重要である。

設計論点には、機能そのものではないが、システム全体に影響するものもある。たとえば、ログ、認証、監査、例外処理、キャッシュ、トランザクション、セキュリティ確認などは、多くの機能に横断的に関わる。こうした論点は\term{横断的関心事}{crosscutting concern} と呼ばれる。

\begin{pointbox}{注意}
横断的関心事を場当たり的に実装すると、同じ処理が各所に重複し、変更時に漏れが起きやすくなる。設計段階で共通方針を決めることが重要である。
\end{pointbox}
